\section{Conclusion}
\label{sec:conclusion}

This paper formalised marginal-carbon-aware deferral of deadline-elastic serverless workloads in the cloud--edge continuum, motivated by two deficits in the carbon-aware shifting literature: the use of the average carbon signal in place of the marginal one, and the use of either a greedy or a long-horizon forecast horizon in place of the short renewable-surplus windows that match how clean hours actually arrive. The model charges each placement the marginal intensity at the time and site of execution, plus an embodied term amortised per served request whenever a warm instance is reused instead of cold-started, exposing the cold-start decision to the carbon budget. The offline optimum knows the full marginal trajectory; the online algorithm does not, and is measured by the competitive ratio against that optimum.

The \textsc{CarbonShift} algorithm is competitive: under adversarial marginal sequences its competitive ratio is $\Theta(\Lambda)$ and the bound is tight up to a constant, and under bounded-variation sequences that match observed grid behaviour the ratio collapses to a small constant that depends on the variation budget and the slack. The algorithm is SLA-safe under a slack threshold that the operator sets, and the slack threshold is the knob that trades carbon savings against SLA risk. The algorithm is non-gameable: a workload that inflates its reported deadline cannot reduce its charged carbon, because the SLA-violation penalty offsets any carbon reduction from a longer window. Three theorems make these claims.

The experimental design pairs public marginal-intensity traces with the real Azure Functions 2019 invocation trace on an open serverless platform, evaluating emissions, SLA compliance, and cold-start cost against four baselines. The results confirm the theory: \textsc{CarbonShift} achieves 1--2\% carbon reduction relative to Greedy while maintaining zero SLA violations, the marginal signal outperforms the average signal by 0.4\%, the horizon U-curve bottoms out at $H=9$ slots as Theorem~\ref{thm:cr-bv} predicts, and the non-gameability guarantee holds with the inflated-deadline charged cost 5.8\% higher than the honest report.

Two directions are open. First, the model assumes capacity-sufficient sites; extending the competitive guarantee to the capacity-bounded regime would make the scheme closer to a production orchestrator. Second, the embodied model assumes a static per-instance budget; extending it to the dynamic case where the budget depends on the instance's future load would close the loop between the reuse decision and the per-request embodied cost more tightly. Both extensions are left for future work.

This paper is submitted to Focus Area~3, Novel Service Frameworks for the Cloud Continuum and Smart Environments, of ICSOC 2026, and addresses the sub-area of Green ICT and Sustainability of Service Systems.
